% ============================================
% DOC-15
% Gerado de molambudos.md — build_miolo.py
% ============================================
\fragtitle{doc}{D}{DOC}{15}{Diário de Oliveira V}
\markboth{DOC-15 --- Diário de Oliveira V}{}

\noindent \textit{A vala existe}

\noindent \textbf{Dr. Heitor Oliveira --- Caderno Pessoal}
\textbf{17 de junho de 1979 --- Fortaleza, Ceará}

\vspace{0.5em}\noindent\rule{\textwidth}{0.5pt}\vspace{0.5em}

\noindent Encontrei.

\noindent A vala existe.

\noindent Passei três dias procurando. Os limites do antigo Campo do Alagadiço, nos arredores de Fortaleza, já não eram legíveis na paisagem. Sem pesquisa especializada, os vestígios eram difíceis de localizar.

\noindent Mas um velho de 90 anos, sentado num bairro de Fortaleza, disse:

\noindent \textit{\textquotedblleft{}A vala? Sei onde é. Fica atrás do que era o curral. Não tem nada lá. Só mato. Mas o mato não cresce direito. A terra é preta. E fedida.\textquotedblright{}}

\noindent Ele me levou.

\noindent Caminhamos 20 minutos por uma estrada de terra. Depois, mato. Depois, uma clareira.

\noindent A terra era preta. Oleosa. O cheiro adocicado subia do chão como vapor.

\noindent \textit{\textquotedblleft{}Aqui\textquotedblright{}\textit{, disse o velho. }\textquotedblleft{}Enterraram os molambudos aqui. Centenas. Milhares. Ninguém sabe quantos.\textquotedblright{}}

\noindent \textit{\textquotedblleft{}Como o senhor sabe?\textquotedblright{}}

\noindent \textit{\textquotedblleft{}Meu pai trabalhou no curral. Ele disse que os soldados jogavam os corpos aqui. De noite. Sem rezar. Sem nada.\textquotedblright{}}

\noindent \textit{\textquotedblleft{}E nunca mais cresceu nada?\textquotedblright{}}

\noindent \textit{\textquotedblleft{}Cresce. Mas cresce torto. Olha.\textquotedblright{}}

\noindent Ele apontou para uma árvore no meio da clareira. O tronco era retorcido, os galhos cresciam para baixo. As folhas eram escuras.

\noindent \textit{\textquotedblleft{}Essa árvore cresceu ali. Bebeu da terra preta. Olha a cor das folhas. Parece... sangue seco.\textquotedblright{}}

\noindent Paguei o velho. Ele foi embora.

\noindent Fiquei sozinho na clareira.

\noindent O cheiro adocicado era forte. Muito forte. Como na cela. Como no meu quarto.

\noindent Sentei no chão. A terra era macia. Fofa. Como se estivesse cheia de buracos.

\noindent Pus a mão na terra preta.

\noindent Ela estava quente.

\noindent Tirei o diário da mochila. Abri na última página.

\noindent A letra estava se formando sozinha:

\noindent \textit{\textquotedblleft{}Você chegou.\textquotedblright{}}

\noindent \textit{\textquotedblleft{}Sim\textquotedblright{}}, eu disse em voz alta.

\noindent \textit{\textquotedblleft{}O que você esperava encontrar?\textquotedblright{}}

\noindent \textit{\textquotedblleft{}A verdade.\textquotedblright{}}

\noindent \textit{\textquotedblleft{}A verdade é que você nunca mais vai embora daqui.\textquotedblright{}}

\noindent Olhei em volta. A clareira. A árvore torta. A terra preta. O cheiro.

\noindent \textit{\textquotedblleft{}Não\textquotedblright{}\textit{, eu disse. }\textquotedblleft{}Eu vou embora. Eu vou voltar para casa. Vou queimar este diário. Vou esquecer.\textquotedblright{}}

\noindent \textit{\textquotedblleft{}Não vai.\textquotedblright{}}

\noindent \textit{\textquotedblleft{}Vou.\textquotedblright{}}

\noindent Fiquei em silêncio.

\noindent A criatura dentro de mim riu.

\noindent \textit{\textquotedblleft{}Você já é meu, Heitor. Desde a primeira página. O ciclo não tem volta.\textquotedblright{}}

\noindent Fiquei sentado na vala até o sol se pôr.

\noindent Quando olhei para o céu, a Rasga Mortalha estava lá. Pousada numa árvore seca. Me olhando.

\noindent Fiquei olhando de volta.

\noindent Até o escuro engolir tudo.

\vspace{0.5em}
\noindent\textit{\textbf{↪ Rotas:} \rota{DOC-16} \rotasep \rota{MEM-04}}